\documentclass[11pt,a4paper]{article}
\usepackage[utf8]{inputenc}
\usepackage[T1]{fontenc}
\usepackage[english]{babel}
\usepackage{geometry}
\geometry{margin=2.5cm}
\usepackage{setspace}
\setstretch{1.05}
\usepackage{longtable}
\usepackage{array}
\usepackage{booktabs}
\usepackage{xcolor}
\usepackage{listings}
\usepackage{hyperref}
\hypersetup{colorlinks=true,linkcolor=blue,urlcolor=blue,citecolor=blue}
\lstdefinestyle{pseudo}{
  basicstyle=\ttfamily\small,
  breaklines=true,
  frame=single,
  columns=fullflexible,
  keepspaces=true,
  showstringspaces=false
}
\title{Meso2D Refactor EN Symbol Reference}
\author{}
\date{}
\begin{document}
\maketitle
\tableofcontents
\newpage
\section*{Meso2D\_refactor\_EN Program Symbol Reference}
\addcontentsline{toc}{section}{Meso2D\_refactor\_EN Program Symbol Reference}

\subsection*{Scope}
\addcontentsline{toc}{subsection}{Scope}
This document lists the classes, methods, and variables used in the main application file Meso2D\_refactor\_EN.py.

The focus is on each symbol's functionality inside the program flow.


\subsection*{Program Classes}
\addcontentsline{toc}{subsection}{Program Classes}

\begin{longtable}{|p{0.475\textwidth}|p{0.475\textwidth}|}
\hline
Class & Purpose in the Program \\\\
\hline
Selector & Initial mode selection window. Lets the user choose circles, ellipses, or polygons before opening the main window. \\\\
\hline
About & Informational label/dialog with project metadata and contact information. \\\\
\hline
mainProgram & Main application window. Handles UI events, project I/O, simulation setup, progress logging, and structure export. \\\\
\hline
Viewer\_image & Secondary dialog used to visualize and save the generated structure image. \\\\
\hline
Controlador & Startup coordinator. Creates the selector window and opens the main window after mode selection. \\\\
\hline
\end{longtable}


\subsection*{External Services and Classes Used}
\addcontentsline{toc}{subsection}{External Services and Classes Used}

\begin{longtable}{|p{0.317\textwidth}|p{0.317\textwidth}|p{0.317\textwidth}|}
\hline
Symbol & Source & Functionality in this Program \\\\
\hline
Ui\_MainWindow & src.ui.Main03\_EN & Generated Qt UI class for the main application window. \\\\
\hline
Menuinicio (Ui\_MainWindow alias) & src.ui.Main\_inicio\_EN & Generated Qt UI class for the mode selector window. \\\\
\hline
Ui\_Dialog\_GV & src.ui.dialog\_GV\_EN & Generated Qt UI class for the image viewer dialog. \\\\
\hline
MiGraphicsView & src.ui.GV\_EN & Custom graphics view widget used to display generated images with viewer controls. \\\\
\hline
SimulationController & src.simulation.orchestration\_EN & Starts and stops the appropriate simulation worker based on selected mode. \\\\
\hline
build\_project\_payload & src.io.project\_EN & Collects current in-memory state into a project dictionary for persistence. \\\\
\hline
save\_project\_to\_file & src.io.project\_EN & Writes serialized project payload to disk. \\\\
\hline
load\_project\_from\_file & src.io.project\_EN & Reads and parses project file from disk. \\\\
\hline
apply\_project\_payload & src.io.project\_EN & Applies loaded payload fields to the active mainProgram instance. \\\\
\hline
export\_structure\_file & src.io.structure\_export\_EN & Performs format-aware structure export for DXF/SVG/JSON/GEO. \\\\
\hline
StructureExportError & src.io.structure\_export\_EN & Domain exception used to report export validation/runtime issues. \\\\
\hline
save\_current\_view & src.visualization.service\_EN & Saves current rendered structure image from the main window context. \\\\
\hline
save\_viewport\_image & src.visualization.service\_EN & Saves rendered image from the viewer dialog context. \\\\
\hline
\end{longtable}


\subsection*{Methods by Class}
\addcontentsline{toc}{subsection}{Methods by Class}

\subsubsection*{Selector}
\addcontentsline{toc}{subsubsection}{Selector}

\begin{longtable}{|p{0.475\textwidth}|p{0.475\textwidth}|}
\hline
Method & Functionality \\\\
\hline
\_\_init\_\_ & Builds selector UI and connects mode buttons to their handlers. \\\\
\hline
pb\_circulos\_pulsar & Emits mode "circulos" and closes selector. \\\\
\hline
pb\_ellipses\_pulsar & Emits mode "ellipses" and closes selector. \\\\
\hline
pb\_poligonos\_pulsar & Emits mode "poligonos" and closes selector. \\\\
\hline
\end{longtable}


\subsubsection*{About}
\addcontentsline{toc}{subsubsection}{About}

\begin{longtable}{|p{0.475\textwidth}|p{0.475\textwidth}|}
\hline
Method & Functionality \\\\
\hline
\_\_init\_\_ & Creates the informational text label and centers text alignment. \\\\
\hline
initUI & Delegates window-centering behavior. \\\\
\hline
center & Moves the widget to screen center based on available geometry. \\\\
\hline
\end{longtable}


\subsubsection*{mainProgram}
\addcontentsline{toc}{subsubsection}{mainProgram}

\begin{longtable}{|p{0.475\textwidth}|p{0.475\textwidth}|}
\hline
Method & Functionality \\\\
\hline
\_\_init\_\_(mode) & Initializes main UI, connects menu/actions and buttons, sets default state, and creates SimulationController. \\\\
\hline
sobre\_programa & Opens About window. \\\\
\hline
close\_application & Asks for confirmation and exits the app if accepted. \\\\
\hline
llenar\_data\_tabla\_dosif & Fills the gradation table from sieve\_size/tpp lists; clears if data is missing. \\\\
\hline
add\_fila & Adds one row to gradation table. \\\\
\hline
open\_project & Handles "open project" workflow with overwrite confirmation when data already exists. \\\\
\hline
open & Reads project file from disk and loads values into UI/state. \\\\
\hline
variables\_limpias & Resets internal state variables to start from a clean project context. \\\\
\hline
new\_project & Starts a new project with optional cleanup confirmation. \\\\
\hline
load\_example & Loads an example project, supporting legacy Spanish filenames. \\\\
\hline
llenar\_data & Reflects internal state values into UI controls and checkboxes. \\\\
\hline
update\_data(diccionario) & Maps legacy key aliases and applies payload values to current object. \\\\
\hline
save\_project & Serializes and writes project data to a text file. \\\\
\hline
plotear(image) & Displays generated image in viewer, updates UI state, and logs simulation completion/time. \\\\
\hline
run & Main entry from Run button: validates replacement of previous structure, updates data, and triggers simulation checks. \\\\
\hline
comprobar\_data\_para\_worker & Verifies required simulation inputs; only starts simulation when all required values are present. \\\\
\hline
simulate & Sets random seed and feature toggles, updates UI progress controls, and delegates execution to SimulationController.start(). \\\\
\hline
stop & Compatibility wrapper for stopping current simulation controller. \\\\
\hline
stop\_simulation & Stops current simulation controller (same behavior as stop). \\\\
\hline
limpiar & Clears generated structure and temporary lists to prepare for another simulation run. \\\\
\hline
update & Pulls values from UI fields into numeric state variables and rebuilds derived areas and gradation arrays. \\\\
\hline
habilitar & Enables/disables related UI fields based on active checkboxes (pores, points, seed, etc.). \\\\
\hline
ver\_structure & Opens image viewer dialog. \\\\
\hline
pasar\_list\_pores(list) & Receives pore list emitted by worker/controller. \\\\
\hline
pasar\_list\_coarse(list) & Receives coarse aggregate list emitted by worker/controller. \\\\
\hline
pasar\_list\_fine(list) & Receives fine aggregate list emitted by worker/controller. \\\\
\hline
pasar\_list\_reactive(list) & Receives reactive point list emitted by worker/controller. \\\\
\hline
exportar\_structure & Generic export entrypoint. Exports to DXF/SVG/JSON/GEO through structure\_export\_EN service. \\\\
\hline
exportar\_structure\_circulos & Legacy/manual DXF export routine for circle-based structure representation. \\\\
\hline
exportar\_structure\_ellipses & Legacy/manual DXF export routine for ellipse-based structures (polyline approximation). \\\\
\hline
points\_elipse (nested helper) & Generates sampled points of a rotated ellipse for DXF polyline export. \\\\
\hline
exportar\_structure\_poligonos & Legacy/manual DXF export routine for polygonal structures, handling Shapely or vertex-list inputs. \\\\
\hline
extract\_polygon\_vertices (nested helper) & Normalizes polygon sources to a clean vertex list for DXF writing. \\\\
\hline
save\_image & Saves current rendered structure image using visualization service. \\\\
\hline
informar2 & Replaces trailing placeholder line (...) in info log or appends a new colored message. \\\\
\hline
log & Appends a colored message to information panel. \\\\
\hline
log\_worker & Appends worker informational message in standard text color. \\\\
\hline
log\_worker\_error & Appends worker error message in error color. \\\\
\hline
show\_progress(percentage) & Updates progress bar value from worker/controller progress signal. \\\\
\hline
\end{longtable}


\subsubsection*{Viewer\_image}
\addcontentsline{toc}{subsubsection}{Viewer\_image}

\begin{longtable}{|p{0.475\textwidth}|p{0.475\textwidth}|}
\hline
Method & Functionality \\\\
\hline
\_\_init\_\_ & Builds dialog UI, creates graphics viewer widget, and wires save/close buttons. \\\\
\hline
save & Saves viewport image through visualization service. \\\\
\hline
close & Hides the image dialog window. \\\\
\hline
\end{longtable}


\subsubsection*{Controlador}
\addcontentsline{toc}{subsubsection}{Controlador}

\begin{longtable}{|p{0.475\textwidth}|p{0.475\textwidth}|}
\hline
Method & Functionality \\\\
\hline
\_\_init\_\_ & Creates selector and binds selected mode signal to main window creation. \\\\
\hline
open\_mainprogram(mode) & Creates and shows mainProgram for selected simulation mode. \\\\
\hline
\end{longtable}


\subsection*{Variables and Attributes}
\addcontentsline{toc}{subsection}{Variables and Attributes}

\subsection*{Module-level Variables}
\addcontentsline{toc}{subsection}{Module-level Variables}

\begin{longtable}{|p{0.475\textwidth}|p{0.475\textwidth}|}
\hline
Variable & Functionality \\\\
\hline
app & QApplication instance, required to run Qt event loop. \\\\
\hline
c & Controlador instance that bootstraps the UI flow. \\\\
\hline
\end{longtable}


\subsection*{Selector Attributes}
\addcontentsline{toc}{subsection}{Selector Attributes}

\begin{longtable}{|p{0.475\textwidth}|p{0.475\textwidth}|}
\hline
Attribute & Functionality \\\\
\hline
mode\_escogido & Qt signal carrying selected mode string to Controlador. \\\\
\hline
\end{longtable}


\subsection*{mainProgram Attributes}
\addcontentsline{toc}{subsection}{mainProgram Attributes}

\subsubsection*{Configuration and Input Data}
\addcontentsline{toc}{subsubsection}{Configuration and Input Data}

\begin{longtable}{|p{0.475\textwidth}|p{0.475\textwidth}|}
\hline
Attribute & Functionality \\\\
\hline
mode & Current simulation mode: circulos, ellipses, or poligonos. \\\\
\hline
sieve\_size & Sieve diameter list used for gradation fractions. \\\\
\hline
tpp & Cumulative passing percentage list aligned with sieve\_size. \\\\
\hline
x, y & Specimen dimensions. \\\\
\hline
Pagg & Aggregate area fraction coefficient. \\\\
\hline
Pporos & Pore area fraction coefficient. \\\\
\hline
dporo\_min, dporo\_max & Minimum and maximum pore diameters. \\\\
\hline
r\_react & Reactive radius parameter (currently not used directly in this file). \\\\
\hline
dpto\_max\_aggregates, dpto\_min\_aggregates & Reactive point diameter limits for aggregates. \\\\
\hline
dpto\_max\_pasta, dpto\_min\_pasta & Reactive point diameter limits for paste matrix. \\\\
\hline
Ppto\_react\_aggregates & Reactive point area fraction in aggregates. \\\\
\hline
Ppto\_react\_pasta & Reactive point area fraction in paste. \\\\
\hline
seed & Random seed used for reproducibility. \\\\
\hline
data & Canonical project key list/payload tracking used for load/save. \\\\
\hline
js & Loaded project dictionary from disk. \\\\
\hline
project & Payload dictionary generated for saving. \\\\
\hline
data\_necesarios & Runtime list used to validate that required values are present before simulation. \\\\
\hline
\end{longtable}


\subsubsection*{Derived Numerical State}
\addcontentsline{toc}{subsubsection}{Derived Numerical State}

\begin{longtable}{|p{0.475\textwidth}|p{0.475\textwidth}|}
\hline
Attribute & Functionality \\\\
\hline
A & Total specimen area (x*y). \\\\
\hline
Aagg & Area allocation per aggregate gradation interval. \\\\
\hline
A\_remanente & Remaining unassigned area after particle generation. \\\\
\hline
A\_pores & Target total pore area. \\\\
\hline
A\_points\_aggregates & Target total reactive-point area over aggregates. \\\\
\hline
A\_points\_pasta & Target total reactive-point area in paste. \\\\
\hline
num\_particulas & Counter used in particle generation workflows. \\\\
\hline
particulas & Aggregate/particle count bookkeeping by fraction. \\\\
\hline
radios & Generic radius list used by worker workflows. \\\\
\hline
radios\_pores & Pore radius list pending placement. \\\\
\hline
radios\_points & Global reactive-point radii trace. \\\\
\hline
radios\_points\_aggregates & Reactive-point radii assigned to aggregate placement. \\\\
\hline
radios\_points\_pasta & Reactive-point radii assigned to paste placement. \\\\
\hline
\end{longtable}


\subsubsection*{Geometry Collections and Export Data}
\addcontentsline{toc}{subsubsection}{Geometry Collections and Export Data}

\begin{longtable}{|p{0.475\textwidth}|p{0.475\textwidth}|}
\hline
Attribute & Functionality \\\\
\hline
list\_aggregates & Runtime aggregate geometry list. \\\\
\hline
todos\_aggregates & Render-ready aggregate patch/object list. \\\\
\hline
list\_aggregates\_coarse & Coarse aggregate export/render list. \\\\
\hline
list\_aggregates\_fine & Fine aggregate export/render list. \\\\
\hline
list\_pores & Pore geometry list used for render/export. \\\\
\hline
todos\_pores & Render-ready pore patch list. \\\\
\hline
list\_ptos\_react & Combined reactive-point list for export/render. \\\\
\hline
list\_ptos\_react\_aggregates & Reactive-point list located on aggregates. \\\\
\hline
list\_ptos\_react\_pasta & Reactive-point list located in paste. \\\\
\hline
todos\_ptos\_react & Render-ready reactive-point patch list. \\\\
\hline
\end{longtable}


\subsubsection*{UI and Execution Control State}
\addcontentsline{toc}{subsubsection}{UI and Execution Control State}

\begin{longtable}{|p{0.475\textwidth}|p{0.475\textwidth}|}
\hline
Attribute & Functionality \\\\
\hline
dlg & Viewer\_image dialog instance. \\\\
\hline
pop & About dialog/label instance. \\\\
\hline
simulation\_controller & Orchestration service that selects/starts/stops the worker by mode. \\\\
\hline
start\_time & Simulation start timestamp used for elapsed-time logging. \\\\
\hline
existe\_structure & Indicates whether a generated structure already exists in current session. \\\\
\hline
todo\_correcto & Indicates whether mandatory data validation passed. \\\\
\hline
check\_pores & Flag to generate pores in current run. \\\\
\hline
check\_points & Flag to generate reactive points in current run. \\\\
\hline
check\_points\_aggregates & Flag to generate reactive points on aggregates. \\\\
\hline
check\_points\_pasta & Flag to generate reactive points in paste. \\\\
\hline
aggregates\_puestos & Flag indicating aggregates were already placed. \\\\
\hline
pores\_puestos & Flag indicating pores were already placed. \\\\
\hline
redColor, blackColor, blueColor & Text colors used for log severity/status visualization. \\\\
\hline
\end{longtable}


\subsection*{Viewer\_image Attributes}
\addcontentsline{toc}{subsection}{Viewer\_image Attributes}

\begin{longtable}{|p{0.475\textwidth}|p{0.475\textwidth}|}
\hline
Attribute & Functionality \\\\
\hline
gv\_viewer & Graphics view widget used to display the rendered structure image. \\\\
\hline
\end{longtable}


\subsection*{Controlador Attributes}
\addcontentsline{toc}{subsection}{Controlador Attributes}

\begin{longtable}{|p{0.475\textwidth}|p{0.475\textwidth}|}
\hline
Attribute & Functionality \\\\
\hline
selector & Selector window instance for initial mode selection. \\\\
\hline
mainprogram & Main window instance created after mode selection. \\\\
\hline
\end{longtable}


\subsection*{UI Attributes Referenced from Ui\_MainWindow}
\addcontentsline{toc}{subsection}{UI Attributes Referenced from Ui\_MainWindow}

These widgets are not created manually in this file, but are heavily used by methods above.


\begin{longtable}{|p{0.475\textwidth}|p{0.475\textwidth}|}
\hline
Attribute & Functionality \\\\
\hline
tabla\_dosif & Gradation table (sieve\_size, tpp). \\\\
\hline
pb\_run, pb\_stop & Start/stop simulation buttons. \\\\
\hline
pb\_add, pb\_remove & Add/remove-update row interactions. \\\\
\hline
pb\_ver\_structure & Opens rendered structure viewer. \\\\
\hline
cb\_pores, cb\_points, cb\_points\_aggregates, cb\_points\_pasta, cb\_seed & Feature toggle checkboxes that control simulation options and field enabling. \\\\
\hline
le\_x, le\_y, le\_seed & Input fields for dimensions and random seed. \\\\
\hline
le\_dmin, le\_dmax, le\_Pporos & Pore parameter input fields. \\\\
\hline
le\_dmin\_ptos\_aggregates, le\_dmax\_ptos\_aggregates, le\_Pptos\_aggregates & Reactive-point-on-aggregates input fields. \\\\
\hline
le\_dmin\_ptos\_pasta, le\_dmax\_ptos\_pasta, le\_Pptos\_pasta & Reactive-point-in-paste input fields. \\\\
\hline
dsb\_coef\_f & Numeric control for aggregate area fraction. \\\\
\hline
progressBar & Simulation progress indicator. \\\\
\hline
tb\_info & Text log panel for status/error/progress messages. \\\\
\hline
a\_Open, a\_GuardarP, a\_New, a\_CargarE, a\_Structure, a\_GuardarI, a\_Exit, a\_About & Menu actions for project and application lifecycle. \\\\
\hline
\end{longtable}


\subsection*{Execution Flow Summary}
\addcontentsline{toc}{subsection}{Execution Flow Summary}

\begin{enumerate}
\item Controlador shows Selector.
\item Selector emits chosen mode.
\item Controlador opens mainProgram(mode).
\item mainProgram gathers inputs, validates required values, and starts SimulationController.
\item Worker emits geometry/image/progress messages back to mainProgram.
\item mainProgram allows visualization and export once structure data is available.
\end{enumerate}

\subsection*{Worker Methodology and Algorithms}
\addcontentsline{toc}{subsection}{Worker Methodology and Algorithms}

This section describes the numerical and geometric methodology used by workers for aggregate and reactive-point generation/placement.


\subsubsection*{Worker Families and Morphology}
\addcontentsline{toc}{subsubsection}{Worker Families and Morphology}

\begin{longtable}{|p{0.475\textwidth}|p{0.475\textwidth}|}
\hline
Worker & Morphology and Data Representation \\\\
\hline
WorkerTodos (src/workers/Worker\_clusters\_EN.py) & Circular aggregates and pores represented as [x, y, r]. Collision checks are mainly vectorized distance checks. \\\\
\hline
WorkerElipses (src/workers/Worker\_ellipses\_EN.py) & Elliptical aggregates represented by [x, y, a, b, angle], converted to Shapely geometry for robust placement collision checks. \\\\
\hline
WorkerPoligonos (src/workers/Worker\_poligonos\_EN.py) & Irregular polygonal aggregates represented as Shapely Polygon objects; pores can be circles or polygons depending on stage/export path. \\\\
\hline
\end{longtable}


\subsubsection*{1) Gradation Preprocessing and Area Budgeting}
\addcontentsline{toc}{subsubsection}{1) Gradation Preprocessing and Area Budgeting}

All workers follow the same base strategy:


\begin{enumerate}
\item Remove ultra-fine classes: gradation\_sin\_extrafinos keeps sieve sizes >= 2 mm.
\item Compute specimen area: A = x * y.
\item Split aggregate area per gradation interval using:
\end{enumerate}
fraccion\_i = (TPP\_i - TPP\_\\{i+1\\}) / (TPP\_0 - TPP\_\\{end\\})

Aagg\_i = fraccion\_i * Pagg * A

\begin{enumerate}
\item Classify each interval into:
\end{enumerate}
\begin{itemize}
\item Coarse fractions: sieve > 4 mm
\item Fine fractions: sieve <= 4 mm
\end{itemize}
\begin{enumerate}
\item Keep A\_remanente between intervals so area not consumed in one step is transferred to the next one.
\end{enumerate}

Functional objective:

\begin{itemize}
\item Preserve global aggregate proportion Pagg while respecting the gradation curve distribution.
\end{itemize}

\subsubsection*{2) Aggregate Size/Shape Sampling by Worker}
\addcontentsline{toc}{subsubsection}{2) Aggregate Size/Shape Sampling by Worker}

\paragraph{WorkerTodos (Circular)}

\begin{itemize}
\item Coarse and fine particles are sampled by random diameter within each sieve interval.
\item Particle area is circle area pi * (d/2)\textasciicircum{}2.
\item New particles are accumulated while interval target area is not exceeded.
\item Output is a radius list for later spatial placement.
\end{itemize}

\paragraph{WorkerElipses (Elliptical)}

\begin{itemize}
\item Diameter d is sampled from sieve interval.
\item Minor semi-axis b = d/2.
\item Aspect ratio is sampled uniformly in [1.0, 1.5], then a = b * aspect.
\item Ellipse area pi * a * b is consumed against interval area budget.
\item Placement later uses real Shapely ellipses generated from a unit circle scaled and rotated.
\end{itemize}

\paragraph{WorkerPoligonos (Polygonal)}

\begin{itemize}
\item For each candidate particle, target area is derived from sampled diameter equivalent area.
\item generar\_poligono\_irregular creates irregular polygons by:
\item Sampling number of sides (5..10)
\item Perturbing radial distances
\item Scaling vertices to match target area
\item Accepting if relative area error < 5\\%
\item Invalid geometries are cleaned using buffer(0).
\end{itemize}

\subsubsection*{3) Spatial Placement of Aggregates}
\addcontentsline{toc}{subsubsection}{3) Spatial Placement of Aggregates}

All workers use stochastic placement with geometric constraints:


\begin{enumerate}
\item Sample random position (and random angle for anisotropic shapes).
\item Enforce domain containment (full shape inside specimen).
\item Reject overlap with already placed aggregates.
\item Reject overlap with already placed pores (if pores were already generated).
\item Repeat up to max attempts (typically 3000) to avoid infinite loops in dense domains.
\end{enumerate}

Differences by worker:


\begin{itemize}
\item WorkerTodos:
\item Distance-only checks with NumPy arrays for circles.
\item Fast and simple, optimized for [x, y, r] datasets.
\end{itemize}

\begin{itemize}
\item WorkerElipses:
\item Shapely intersects/contains checks on true ellipse geometries.
\item More robust for anisotropic shapes and rotations.
\end{itemize}

\begin{itemize}
\item WorkerPoligonos:
\item Shapely polygon intersection/containment checks.
\item Includes geometry cleaning and explicit placement loop bounds.
\end{itemize}

\subsubsection*{4) Pore Computation and Placement}
\addcontentsline{toc}{subsubsection}{4) Pore Computation and Placement}

Pore methodology is also budget-based:


\begin{enumerate}
\item Target pore area = A * Pporos.
\item Sample pore radius in [dporo\_min/2, dporo\_max/2].
\item Subtract each pore area from remaining target until threshold.
\item Place pores randomly while enforcing:
\end{enumerate}
\begin{itemize}
\item Full containment in specimen
\item No overlap with aggregates
\item No overlap with previously placed pores
\end{itemize}

WorkerPoligonos and WorkerElipses use geometry-aware checks when aggregate geometry is non-circular.


\subsubsection*{5) Reactive Point Computation}
\addcontentsline{toc}{subsubsection}{5) Reactive Point Computation}

Reactive points are generated in two channels:


\begin{itemize}
\item On aggregates (Ppto\_react\_aggregates, dpto\_min\_aggregates, dpto\_max\_aggregates)
\item In paste (Ppto\_react\_pasta, dpto\_min\_pasta, dpto\_max\_pasta)
\end{itemize}

Both channels use area budgeting:


\begin{enumerate}
\item Target reactive-point area = A * ratio.
\item Sample radius in diameter range.
\item Subtract circle area pi * r\textasciicircum{}2 until budget is consumed.
\item Keep separate radius lists for aggregate channel and paste channel.
\end{enumerate}

\subsubsection*{6) Reactive Point Placement: Clusters and Noise}
\addcontentsline{toc}{subsubsection}{6) Reactive Point Placement: Clusters and Noise}

Workers implement two distinct spatial generators for reactive points:


\paragraph{A) Cluster-based placement (typically for points on aggregates)}

\begin{itemize}
\item Generate cluster centers in the domain.
\item Number of clusters is adapted from domain size using influence area:
\item radio\_influencia = 3 * sigma
\item area\_cluster = pi * radio\_influencia\textasciicircum{}2
\item num\_clusters = max(k\_minimo, ceil(domain\_area / area\_cluster))
\item Candidate points are sampled from normal distributions around random centers:
\end{itemize}
x \textasciitilde{} N(cx, sigma), y \textasciitilde{} N(cy, sigma)

\begin{itemize}
\item Candidates must satisfy placement constraints:
\item Inside domain
\item No overlap with pores and existing points
\item Intersect aggregates when channel is "on aggregates"
\end{itemize}

Effect:

\begin{itemize}
\item Creates spatial grouping and localized concentration of reactive points.
\end{itemize}

\paragraph{B) Simplex-noise placement (typically for points in paste)}

\begin{itemize}
\item Uses OpenSimplex noise field with fixed seed (42 in current code).
\item Candidate point is accepted by noise gate:
\end{itemize}
valor = (noise2(x*escala, y*escala) + 1)/2

accept if valor >= umbral

\begin{itemize}
\item Then standard geometric constraints are checked.
\end{itemize}

Main controls:

\begin{itemize}
\item escala: spatial frequency of noise map.
\item umbral: selectivity (higher threshold -> fewer accepted zones).
\end{itemize}

Effect:

\begin{itemize}
\item Produces non-uniform, organic spatial patterns without sharp grid artifacts.
\end{itemize}

\subsubsection*{7) Collision Logic Summary}
\addcontentsline{toc}{subsubsection}{7) Collision Logic Summary}

\begin{longtable}{|p{0.475\textwidth}|p{0.475\textwidth}|}
\hline
Context & Acceptance Rule \\\\
\hline
New aggregate & Must be fully inside domain and non-overlapping with existing aggregates (and pores if present). \\\\
\hline
New pore & Must be inside domain and non-overlapping with aggregates and pores. \\\\
\hline
New reactive point on aggregate & Must intersect at least one aggregate and avoid forbidden overlaps. \\\\
\hline
New reactive point in paste & Must avoid aggregate overlap (except intended logic per worker), pass noise/cluster sampling gate, and avoid forbidden overlaps. \\\\
\hline
\end{longtable}


\subsubsection*{8) Numerical/Implementation Safeguards}
\addcontentsline{toc}{subsubsection}{8) Numerical/Implementation Safeguards}

\begin{itemize}
\item Explicit max attempt loops to avoid deadlocks in saturated domains.
\item Geometry repair (buffer(0)) for invalid polygons/derived shapes.
\item Progress signaling by pipeline stage.
\item Stop checks (check\_stop) interleaved in loops for responsive cancellation.
\item Emission ordering in geometry-based workers: data lists first, image later, to avoid stale export payloads.
\end{itemize}

\subsubsection*{9) Worker Pipeline (Common Pattern)}
\addcontentsline{toc}{subsubsection}{9) Worker Pipeline (Common Pattern)}

\begin{enumerate}
\item Gradation cleanup.
\item Area split by fraction.
\item Coarse aggregate generation by area.
\item Fine aggregate generation by area.
\item Aggregate placement in domain.
\item Optional pore generation and placement.
\item Optional reactive-point generation and placement (clusters/noise as configured).
\item Render/emit geometry and image to UI.
\end{enumerate}

\subsection*{Pseudocode by Stage}
\addcontentsline{toc}{subsection}{Pseudocode by Stage}

This pseudocode summarizes the implementation logic in a code-maintenance-friendly format.


\subsubsection*{A) Common Simulation Pipeline}
\addcontentsline{toc}{subsubsection}{A) Common Simulation Pipeline}

\begin{lstlisting}[style=pseudo]
function simulate_worker(params, flags):
	assert required inputs are valid

	gradation = remove_ultra_fines(params.sieve_size, threshold=2.0)
	area_budget = split_area_by_gradation(gradation, params.tpp, params.Pagg, params.x * params.y)

	coarse_budget, fine_budget = classify_by_sieve(area_budget, split_mm=4.0)

	coarse_shapes = generate_coarse_aggregates(coarse_budget)
	fine_shapes = generate_fine_aggregates(fine_budget)

	placed_coarse, placed_fine = place_aggregates(coarse_shapes, fine_shapes, domain)

	if flags.check_pores:
		pore_radii = generate_pore_radii(params.Pporos, params.dporo_min, params.dporo_max)
		placed_pores = place_pores(pore_radii, domain, placed_coarse + placed_fine)

	if flags.check_points:
		r_agg = generate_reactive_radii(params.Ppto_react_aggregates, params.dpto_min_aggregates, params.dpto_max_aggregates)
		r_paste = generate_reactive_radii(params.Ppto_react_pasta, params.dpto_min_pasta, params.dpto_max_pasta)

		placed_points_agg = place_points_on_aggregates_clustered(r_agg, domain, aggregates, pores)
		placed_points_paste = place_points_in_paste_simplex(r_paste, domain, aggregates, pores)

	emit geometry lists first
	emit rendered image second
\end{lstlisting}


\subsubsection*{B) Gradation and Area Split}
\addcontentsline{toc}{subsubsection}{B) Gradation and Area Split}

\begin{lstlisting}[style=pseudo]
function split_area_by_gradation(sieve, tpp, Pagg, A_total):
	for i in 0 .. len(sieve)-2:
		fraction = (tpp[i] - tpp[i+1]) / (tpp[0] - tpp[-1])
		Aagg_i = fraction * Pagg * A_total
		store Aagg_i
	return Aagg
\end{lstlisting}


\begin{lstlisting}[style=pseudo]
function consume_budget_with_remainder(interval_budget, remaining):
	target = interval_budget + remaining
	consumed = 0
	while target - consumed > epsilon:
		candidate = sample_particle_area()
		if consumed + candidate < target:
			consumed += candidate
			accept particle
		else:
			break
	return accepted_particles, target - consumed
\end{lstlisting}


\subsubsection*{C) Aggregate Generation per Worker}
\addcontentsline{toc}{subsubsection}{C) Aggregate Generation per Worker}

\paragraph{C1) Circular Worker (WorkerTodos)}

\begin{lstlisting}[style=pseudo]
function sample_circle_from_sieve_interval(d_min, d_max):
	d = uniform(d_min, d_max)
	r = d / 2
	area = pi * r^2
	return r, area
\end{lstlisting}


\paragraph{C2) Elliptical Worker (WorkerElipses)}

\begin{lstlisting}[style=pseudo]
function sample_ellipse_from_sieve_interval(d_min, d_max):
	d = uniform(d_min, d_max)
	b = d / 2
	aspect = uniform(1.0, 1.5)
	a = b * aspect
	area = pi * a * b
	return (a, b), area
\end{lstlisting}


\begin{lstlisting}[style=pseudo]
function build_shapely_ellipse(x, y, a, b, angle):
	unit_circle = Point(x, y).buffer(1)
	ellipse = scale(unit_circle, a, b)
	ellipse = rotate(ellipse, angle, origin=(x, y))
	return ellipse
\end{lstlisting}


\paragraph{C3) Polygonal Worker (WorkerPoligonos)}

\begin{lstlisting}[style=pseudo]
function generate_irregular_polygon(target_area, max_tries=100):
	repeat max_tries:
		n_sides = random_int(5, 10)
		angles = equally_spaced(0, 2*pi, n_sides)
		radial_noise = uniform(0.8, 1.2, n_sides)

		vertices = polar_to_cartesian(radial_noise, angles)
		center vertices around origin

		current_area = polygon_area(vertices)
		if current_area == 0:
			continue

		scale_factor = sqrt(target_area / current_area)
		vertices = vertices * scale_factor

		if relative_area_error(vertices, target_area) < 0.05:
			polygon = Polygon(vertices)
			if invalid: polygon = polygon.buffer(0)
			return polygon
	return None
\end{lstlisting}


\subsubsection*{D) Aggregate Placement Core}
\addcontentsline{toc}{subsubsection}{D) Aggregate Placement Core}

\begin{lstlisting}[style=pseudo]
function place_shape_list(shapes, domain, placed_aggregates, placed_pores, max_attempts=3000):
	for shape in shapes:
		placed = false
		for attempt in 1..max_attempts:
			candidate = random_translate_and_optional_rotate(shape)

			if not domain.contains(candidate):
				continue
			if intersects_any(candidate, placed_aggregates):
				continue
			if intersects_any(candidate, placed_pores):
				continue

			placed_aggregates.append(candidate)
			placed = true
			break

		if not placed:
			discard shape to avoid infinite loops in dense domains
\end{lstlisting}


\subsubsection*{E) Pore Radius Generation and Placement}
\addcontentsline{toc}{subsubsection}{E) Pore Radius Generation and Placement}

\begin{lstlisting}[style=pseudo]
function generate_pore_radii(A_total, Pporos, dmin, dmax):
	A_target = A_total * Pporos
	radii = []
	while A_target > pi * (dmin/2)^2:
		r = uniform(dmin/2, dmax/2)
		radii.append(r)
		A_target -= pi * r^2
	return radii
\end{lstlisting}


\begin{lstlisting}[style=pseudo]
function place_pores(radii, domain, placed_aggregates, placed_pores):
	for r in radii:
		try random positions until accepted:
			fully inside domain
			no overlap with aggregates
			no overlap with previous pores
\end{lstlisting}


\subsubsection*{F) Reactive Point Radius Generation}
\addcontentsline{toc}{subsubsection}{F) Reactive Point Radius Generation}

\begin{lstlisting}[style=pseudo]
function generate_reactive_radii(A_total, ratio, dmin, dmax):
	A_target = A_total * ratio
	radii = []
	while A_target > pi * (dmin/2)^2:
		r = uniform(dmin/2, dmax/2)
		radii.append(r)
		A_target -= pi * r^2
	return radii
\end{lstlisting}


\subsubsection*{G) Cluster-Based Reactive Point Placement}
\addcontentsline{toc}{subsubsection}{G) Cluster-Based Reactive Point Placement}

\begin{lstlisting}[style=pseudo]
function generate_cluster_centers(domain, k_min, sigma):
	influence_radius = 3 * sigma
	cluster_area = pi * influence_radius^2
	k = max(k_min, ceil(domain.area / cluster_area))
	return [uniform_point_in_domain() for _ in 1..k]
\end{lstlisting}


\begin{lstlisting}[style=pseudo]
function place_points_clustered(radii, centers, sigma, constraints):
	for r in radii:
		placed = false
		repeat up to max_attempts:
			(cx, cy) = random_choice(centers)
			x = normal(cx, sigma)
			y = normal(cy, sigma)

			if constraints_are_satisfied(x, y, r):
				save point
				placed = true
				break
		if not placed:
			discard radius
\end{lstlisting}


\subsubsection*{H) Simplex-Noise Reactive Point Placement}
\addcontentsline{toc}{subsubsection}{H) Simplex-Noise Reactive Point Placement}

\begin{lstlisting}[style=pseudo]
function place_points_simplex(radii, simplex_seed=42, scale=0.05, threshold=0.3, constraints):
	noise = OpenSimplex(seed=simplex_seed)

	for r in radii:
		placed = false
		repeat up to max_attempts:
			x, y = uniform_point_in_domain()
			v = (noise2(x * scale, y * scale) + 1) / 2

			if v < threshold:
				continue
			if constraints_are_satisfied(x, y, r):
				save point
				placed = true
				break
		if not placed:
			discard radius
\end{lstlisting}


\subsubsection*{I) Constraint Evaluation Logic}
\addcontentsline{toc}{subsubsection}{I) Constraint Evaluation Logic}

\begin{lstlisting}[style=pseudo]
function constraints_are_satisfied(x, y, r):
	if not fully_inside_domain(x, y, r):
		return false
	if overlaps_existing_points(x, y, r):
		return false
	if overlaps_pores(x, y, r):
		return false

	if channel == "on_aggregates":
		if not intersects_at_least_one_aggregate(x, y, r):
			return false
	else if channel == "paste":
		if overlaps_forbidden_aggregate_region(x, y, r):
			return false

	return true
\end{lstlisting}


\subsubsection*{J) Practical Tuning Guide}
\addcontentsline{toc}{subsubsection}{J) Practical Tuning Guide}

\begin{lstlisting}[style=pseudo]
If point distribution is too uniform:
	increase sigma in cluster mode
	reduce simplex threshold

If too many placement failures:
	reduce target ratios (Pagg, Pporos, reactive ratios)
	reduce min diameters
	increase max attempts carefully

If run is too slow:
	simplify geometry checks (circle mode)
	lower sampling resolution for shape approximations
	reduce plotting/export detail during development
\end{lstlisting}


\end{document}




